\documentclass[10pt,a4paper]{amsart}
\usepackage[margin=19mm]{geometry}
\usepackage{amsmath,amssymb,mathtools}
\usepackage[scr=rsfs,cal=euler]{mathalfa}
\usepackage{xfrac}
\usepackage[unicode,hidelinks,bookmarks=false]{hyperref}
\newcommand{\ie}{i.e.\ }
\newcommand{\cf}{cf.\ }
\newcommand{\resp}{respectively\ }
\newcommand{\Subsection}{\subsection}
\providecommand{\nobreakdash}{\nobreak\hbox{-}}
\setlength{\emergencystretch}{2em}
\multlinegap0pt
\usepackage{bm}
\usepackage{bbm}
\usepackage{dsfont}
\usepackage{xspace}
\usepackage{cancel}
\usepackage{mathtools}
\usepackage{amsthm}
\makeatletter
\@ifundefined{theorem}{\newtheorem{theorem}{Theorem}}{}
\@ifundefined{proposition}{\newtheorem{proposition}[theorem]{Proposition}}{}
\makeatother
\DeclareMathOperator{\E}{\mathcal{E}}
\DeclareMathOperator{\Fc}{\mathcal{F}}
\DeclareMathOperator{\Oc}{\mathcal{O}}
\DeclareMathOperator{\Sym}{Sym}
\title[Positivity, plethysm and hyperbolicity of Siegel varieties i]{Positivity, plethysm and hyperbolicity of Siegel varieties in positive characteristic}
\author{Thibault Alexandre}
\date{Source: arXiv:2206.05804v3 (CC BY 4.0)}
\begin{document}
\maketitle

\noindent\emph{Source and licence.} Positivity, plethysm and hyperbolicity of Siegel varieties in positive characteristic. Version arXiv:2206.05804v3 (CC BY 4.0), distributed under the Creative Commons Attribution 4.0 International licence, as recorded in the arXiv licence metadata for this exact version and shown on the abstract page https://arxiv.org/abs/2206.05804v3. Full rights notice: licenses/arxiv-2206.05804-CC-BY-4.0.txt.\par\smallskip

\noindent\textit{Editorial scope.} Counted unit: the Proposition whose source label is \texttt{nD} in the article (source0.tex lines 1068--1074, source label \texttt{nD}), delivered together with the article's own complete original proof of it (source0.tex lines 1075--1139, one \texttt{proof} environment of 65 lines). No mathematics is written, repaired or re-derived: statement and proof are the article's own text line for line. Required context retained uncounted: amp\_criterion (lines 850--862). Every definition, display and label of those lines is retained unchanged. Omitted from this package and not redistributed: 1--849, 863--1067, 1140--2322. Nothing inside a retained excerpt is omitted or altered: every retained body is delivered exactly as the article's lines are written, so no omission receipt of any kind accompanies this package. Citations: the retained excerpts carry no citation command at all, so no bibliography entry of the article is transcribed and no reference is redistributed. Reference adaptation: none was needed and none was made; every reference inside the delivered unit resolves inside it, so no unresolved reference or citation is produced. Printed slips kept literal: no printed word, symbol, hypothesis or number of the counted statement or of its proof was altered. Layout and class: the article's own class is \texttt{amsart} with packages including inputenc, amsmath, amsfonts, amssymb, amsthm, mathtools, geometry, tikz-cd, enumitem, hyperref, placeins, ytableau, verbatim, pgfplots; the wrapper is the lane's pinned \texttt{amsart} editorial wrapper at 10pt on A4, which restates the theorem-like environments the retained text needs and adds the article's own macro definitions that the retained text needs (\texttt{\textbackslash E}, \texttt{\textbackslash Fc}, \texttt{\textbackslash Oc}, \texttt{\textbackslash Sym}), with extra wrapper packages (bm, bbm, dsfont, xspace, cancel, mathtools). The delivered numbering is the unit's own. Assets: no retained span contains a figure environment, a graphics-inclusion command, a PDF-page inclusion, a table or a TikZ picture; no image file is copied into this package, the package declares no asset dependency and no image file is redistributed with it. Licence: version arXiv:2206.05804v3 of this article is distributed under the Creative Commons Attribution 4.0 International licence, as recorded in the arXiv licence metadata for this exact version and shown on the abstract page https://arxiv.org/abs/2206.05804v3. A full rights notice is delivered at licenses/arxiv-2206.05804-CC-BY-4.0.txt. Characters: every retained excerpt is pure ASCII, so the delivered engine, XeLaTeX with its default Latin Modern text font, reports no missing character.
\begin{proposition}\label{amp_criterion}
Let $\E$ be a vector bundle on $X$. The following assertions are equivalent.
\begin{enumerate}
\item $\E$ is ample on $X$.
\item For all coherent sheaf $\Fc$ on $X$, there is an integer $n_0$ such that $\Fc \otimes \Sym^{n}\E$ is globally generated for all $n \geq n_0$.
\item For all coherent sheaf $\Fc$ on $X$, there is an integer $n_0$ such that the cohomology groups $H^i(X,\Fc \otimes \Sym^{n}\E)$ vanishes for all $i>0$, $n \geq n_0$.
\item There exists a real number $\varepsilon > 0$ such that for all finite morphism $g : C \rightarrow X$ where $C$ is a curve, we have
\begin{equation*}
\delta(g^*\E) \geq \varepsilon \lVert g_*C \rVert.
\end{equation*}
Recall that $\delta(g^*\E)$ is the minimum of the degrees of quotient line bundles of $g^*\E$ and $\lVert \cdot \rVert$ denotes a norm on $A_1(X)$, the $k$-vector space of $1$-cycles modulo linear equivalence.
\end{enumerate}
\end{proposition}

\begin{proposition}\label{nD}
Let $\E$ be a vector bundle on $X$ and $n\geq 1$ an integer. The following assertion are equivalent.
\begin{enumerate}
\item $\E$ is $(\varphi,D)$-ample
\item $\E$ is $(\varphi,nD)$-ample
\end{enumerate}
\end{proposition}

\begin{proof}
Assume that $\E$ is $(\varphi,D)$-ample and consider $r_0 \geq 1$ such that $$\E^{(p^r)}(-D)$$ is ample for all $r \geq r_0$. By Barton's numerical criterion of ampleness recalled in assertion $(4)$ of proposition \ref{amp_criterion}, for all $r \geq r_0$, we have a real number $\varepsilon_r > 0$ such that for all finite morphism $g : C \rightarrow X$ where $C$ is a smooth projective curve over $k$, we have
\begin{equation*}
\delta(g^*\E^{(p^r)}(-D)) \geq \varepsilon_r\lVert g_*C \rVert
\end{equation*} 
which is equivalent to
\begin{equation*}
\delta(g^*\E^{(p^r)}) -D\cdot C \geq \varepsilon_r\lVert g_*C \rVert
\end{equation*}
where $D\cdot C$ is the degree of the line bundle $g^*\Oc_X(D) = \Oc_X(D)_{|C}$ (it is also equal to the intersection number of $D$ with $C$). If $D\cdot C\leq 0$, then 
\begin{equation*}
\begin{aligned}
\delta(g^*\E^{(p^r)}(-nD)) &= \delta(g^*\E^{(p^r)})-D\cdot C - (n-1)D\cdot C \\
&\geq \delta(g^*\E^{(p^r)})-D\cdot C \\
&\geq \varepsilon_r\lVert g_*C \rVert
\end{aligned}
\end{equation*}
for all $r\geq r_0$. If $D\cdot C > 0$, we take $r_1 \geq r_0$ such that $r_1 \geq r_0 + \log_p(n)$ and let $r\geq r_1$ be an integer. Since $\E^{(p^r)}(-D)$ is ample, the bundle 
\begin{equation*}
(\E^{(p^{r_0})}(-D))^{(p^{r-r_0})} = \E^{(p^r)}(-p^{r-r_0}D)
\end{equation*}
is ample and we have
\begin{equation*}
\delta(g^*\E^{(p^r)}(-p^{r-r_0}D)) \geq \varepsilon_r^\prime \lVert g_*C \rVert
\end{equation*}
for some real number $\varepsilon_r^\prime > 0$. Thus,
\begin{equation*}
\begin{aligned}
\delta(g^*\E^{(p^r)}(-nD)) &= \delta(g^*\E^{(p^r)}(-p^{r-r_0}D)) +(p^{r-r_0}-n)D\cdot C \\
&\geq  \delta(g^*\E^{(p^r)}(-p^{r-r_0}D)) \\
&\geq \varepsilon_r^\prime \lVert g_*C \rVert.
\end{aligned}
\end{equation*}
In conclusion, we have
\begin{equation*}
\delta(g^*\E^{(p^r)}(-nD)) \geq \min(\varepsilon_r,\varepsilon_r^\prime) \lVert g_*C \rVert
\end{equation*}
for all $r \geq r_1$ and all $g : C \rightarrow X$, which means that $\E$ is $(\varphi,nD)$-ample. Inversely, consider an integer $r_0 \geq 1$ such that for all $r \geq r_0$, we have a real number $\varepsilon_r > 0$ such that for all finite morphism $g : C \rightarrow X$ where $C$ is a smooth projective curve over $k$, we have
\begin{equation*}
\delta(g^*\E^{(p^r)}) -nD\cdot C \geq \varepsilon_r\lVert g_*C \rVert.
\end{equation*}
If $D\cdot C\geq 0$, we have 
\begin{equation*}
\begin{aligned}
\delta(g^*\E^{(p^r)}(-D)) &= \delta(g^*\E^{(p^r)}(-nD)) +(n-1)D\cdot C \\
&\geq \delta(g^*\E^{(p^r)}(-nD)) \\
&\geq \varepsilon_r\lVert g_*C \rVert
\end{aligned}
\end{equation*}
for all $r\geq r_0$.
Consider an integer $r_1 \geq r_0 + \log_pn$. If $D\cdot C< 0$, we have 
\begin{equation*}
\begin{aligned}
p^{r-r_0}\delta(g^*\E^{(p^{r_0})}(-D)) &\geq \delta(g^*\E^{(p^r)}(-p^{r-r_0}D)) \\
&\geq \delta(g^*\E^{(p^r)}(-nD)) + (n-p^{r-r_0})D\cdot C \\
&\geq \delta(g^*\E^{(p^r)}(-nD)) \\
&\geq \varepsilon_r\lVert g_*C \rVert.
\end{aligned}
\end{equation*}
for all $r \geq r_1$. In conclusion, we have
\begin{equation*}
\delta(g^*\E^{(p^r)}(-D)) \geq \frac{\varepsilon_r}{p^{r-r_0}} \lVert g_*C \rVert
\end{equation*}
for all $r \geq r_1$ and all $g : C \rightarrow X$, which means that $\E$ is $(\varphi,D)$-ample.
\end{proof}

\end{document}
